\def\stat{korchazhkina}





\def\tit{ОПТИМИЗАЦИЯ ПОИСКА ПРИ~РЕШЕНИИ ПЕРЕБОРНЫХ ЗАДАЧ В~УГЛУБЛЕННОМ КУРСЕ 


ИНФОРМАТИКИ НА~УРОВНЕ ОСНОВНОГО ОБЩЕГО ОБРАЗОВАНИЯ}





\def\titkol{Оптимизация поиска при~решении переборных задач в~углубленном курсе 


информатики} % на~уровне основного общего образования}





\def\aut{О.\,М.~Корчажкина$^1$}





\def\autkol{О.\,М.~Корчажкина}





\titel{\tit}{\aut}{\autkol}{\titkol}





\index{Корчажкина О.\,М.}


\index{Korchazhkina O.\,M.}





%{\renewcommand{\thefootnote}{\fnsymbol{footnote}}


%\footnotetext[1]


%{Работа выполнена в~МСЦ РАН в~рамках государственного задания по теме FNEF-2022-0014.}} 





\renewcommand{\thefootnote}{\arabic{footnote}}


\footnotetext[1]{Федеральный исследовательский центр <<Информатика и~управ\-ле\-ние>> Российской академии наук, 


\mbox{olgakomax@gmail.com}}





 \label{st\stat}





 \thispagestyle{headings}


 


 \vspace*{-10pt} 








     \Abst{Рассматриваются способы активизации логического мышления учащихся 


с~целью формирования алгоритмических навыков на уроках информатики в~средней школе. 


Анализируются методы оптимизации в~традиционных криптоарифметических задачах, 


составляющих подкласс переборных задач, предложенных Гербертом Саймоном для 


иллюстрации способов мышления, когда требуется осуществлять многочисленные переборы 


возможных комбинаций. Предложен визуальный метод оптимизации поиска применительно к~решению двух задач Саймона, основанный на закономерностях переноса разрядов при 


суммировании и~представленный в~виде простых алгебраических соотношений. Метод легко 


поддается алгоритмизации и~может быть реализован с~помощью технологии динамического 


программирования, используемой для оптимизации решения переборных задач.}


     


     \KW{оптимизация поиска; переборная задача; криптоарифметическая задача; 


динамическое программирование}





\DOI{10.14357\!/\!08696527220414}  





\vskip 10pt plus 9pt minus 6pt





 \thispagestyle{headings}


 


 \vspace*{-16pt}





\section{Введение}





\vspace*{-2pt}





    Решение переборных задач может служить эффективным инструментом 


развития логического мышления как необходимого компонента инженерной 


культуры учащихся старших классов средней школы~[1]. Переборные  


задачи~--- это задачи, в~которых требуется найти оптимальное решение среди 


конечного, но часто весьма значительного набора возможных альтернатив 


путем сравнения их между собой по заданным критериям или значениям 


параметров. Основной целью при этом выступает умение строить алгоритмы 


перебора комбинаторных объектов (последовательностей, перестановок, 


подмножеств, групп) в~различных комбинациях~[2]. 


   % 


    При этом важно сначала представить решение переборной задачи <<на 


бумаге>>~--- в~виде некоторого алгоритма, который затем по причине 


сложности и~значительного объема <<ручных>> вычислительных процедур 


реализовать с~помощью программных средств.





 Наиболее эффективно~--- 


рационально, быстро и~красиво~--- решение переборных задач воплощается 


средствами динамического программирования~[3--6]. Это объясняется тем, что 


технология динамического программирования эффективно работает для 


реализации характерных для переборных задач рекурсивных алгоритмов, 


которые выдают одинаковые результаты для одинаковых подзадач. Поэтому 


нет необходимости проводить повторные вычисления на каждом следующем 


шаге: достаточно сохранить промежуточные результаты в~соответствующих 


архивах данных, что экономит время, затрачиваемое на вычисление конечного 


результата~\cite[с.~49]{3-kor}. 


    


    Отсюда следует, что формирование и~развитие способности к~составлению 


алгоритмов оптимального поиска при решении переборных задач 


<<вручную>>~--- это первый обязательный (пропедевтический) шаг 


к~обучению конкретным методам динамического программирования, владение 


которыми укрепляет основы логического, алгоритмического и~математического 


мышления учащихся.


    


\section{Оптимизация поиска в~двух криптоарифметических задачах 


Герберта Саймона}





    Методология поиска оптимальных решений, разработанная выдающимся 


американским математиком, экономистом, психологом и~кибернетиком 


Гербертом Саймоном (1916--2001), базируется на логических основаниях 


рационального выбора среди заданных альтернатив с~помощью совокупности 


методов, <<позволяющих на практике отыскивать среди имеющихся 


альтернатив оптимальную>>~\cite[с.~78--94]{7-kor}. Наблюдаемые при 


решении некоторых задач с~помощью созданного им в~1957~г.\ совместно 


с~Алленом Ньюэллом универсального решателя задач (\textit{General Problem 


Solver}~--- GPS) возможные случаи <<комбинаторного взрыва>> (см.\ примеры 


в~\cite[с.~61--93]{8-kor}) заставили ученого вплотную заняться методами 


оптимизации поиска, которые могли бы предотвращать подобные коллапсы как 


при работе компьютерных систем, так и~при переборе решений в~<<ручном>> 


режиме. 


    


    Перебор всех возможных альтернатив, к~которому прибегают в~случае 


неумения оптимизировать (так называемый <<лобовой>> поиск), сводится 


к~рассмотрению максимально возможного числа вариантов решения задачи 


и~сравнению каждого варианта с~заданными условиями. Оптимальный поиск изначально предполагает обратное движение~--- от 


условия задачи к~возможным альтернативам, когда сжатие области поиска 


происходит уже на этапе первоначального отсечения заведомо проигрышных 


вариантов, очевидно не удовлетворяющих условию задачи. Такой метод 


оптимизации Саймон называет минимаксным. 


   % 


    Для этого полезно иметь перед глазами полную (насколько это возможно) 


картину наиболее перспективных решений, составленную на основе 


неоспоримых условий или представленную в~виде математических выражений, 


нарушение которых приведет к~заведомо тупиковому результату~--- если речь 


идет о задачах, поддающихся алгоритмизации и~вычислению. Поэтому 


оптимальный поиск для получения верного решения переборной задачи, 


требующей, как правило, значительных усилий, за\-тра\-чи\-ва\-емых вручную, 


сводится к~<<допоисковому>> ограничению числа рассматриваемых 


альтернатив на основе правил и~закономерностей, диктуемых условиями задачи 


или выводимых из задачной ситуации. Затем на каждом дальнейшем шаге 


(итерации) после сравнения промежуточного результата с~условиями задачи 


(в~том числе и~с~изменяющимися условиями) непременно происходит 


дальнейшее сужение области поиска. Герберт Саймон назвал эти изменения 


<<процессом сбора информации о структуре задачи, которая в~конечном счете 


позволит найти решение>>~\cite[с.~90]{7-kor}. 


    


    К криптоарифметическим задачам относятся переборные задачи, где 


в~словах нужно заменить буквы цифрами от~0 до~9. Чаще всего слова 


записываются в~столбик так, чтобы сумма первых двух слов в~цифровом 


выражении составила число, полученное при замене цифрами букв в~третьем 


слове. Например, две классические криптоарифметические задачи, 


предложенные Гербертом Саймоном, выглядят следующим образом: 


    \begin{center}


    \tabcolsep=2pt


    \begin{tabular}{clcl}


              &\textit{Задача 1} & &\textit{Задача 2}\\


                  &\textbf{DONALD}& &\hphantom{\textbf{D}}\textbf{C\,R\,O\,S\,S}\\[-6pt]


              $+$ &&$+$&\\[-6pt]


              &\underline{\textbf{GERALD}}   &&   \underline{\hphantom{\textbf{D}}\textbf{R\,O\,A\,D\,S}}\\


              &\textbf{ROBERT}, где 


\textbf{D}\;=\;5;&&\textbf{D\,A\,N\,G\,E\,R}.\\


              \end{tabular}


              \end{center}


 


 В~\cite[с.~39--43]{7-kor} автор предлагает рассмотреть четыре стратегии 


поиска при решении задачи~1. Кратко опишем каждую из них.


 


    \smallskip


    


    \textbf{Стратегия~1} (<<Лобовой>> поиск). Механически перебрать все 


возможные варианты замены десяти букв десятью цифрами. Таких переборов 


будет~10!, или 3\,628\,800.


    


    \smallskip


    


    \textbf{Стратегия 2} (Ограниченный <<лобовой>> поиск). Начинаем 


справа подыскивать подстановки для букв \textbf{D}, \textbf{T}, \textbf{L}, \textbf{R}, \textbf{A}, \textbf{E}, \textbf{N}, \textbf{B}, \textbf{O} и~\textbf{G} 


в~естественном порядке: 1, 2, 3, 4, 5, 6, 7, 8, 9, 0. Известно заранее, что 


$\mathbf{D} \hm= 5$, поэтому будем находить соответствие между рядом из 


девяти букв \textbf{T}, \textbf{L}, \textbf{R}, \textbf{A}, \textbf{E}, \textbf{N}, \textbf{B}, \textbf{O} и~\textbf{G} 


и~рядом из девяти цифр~1, 2, 3, 4, 6, 7, 8, 9 и~0. 


    


    Поскольку в~1-м разряде $2\mathbf{D}\hm = \mathbf{T}$, то $\mathbf{T}


    \hm = 0$, а~единица переносится во 2-й разряд. Далее ряды букв и~цифр 


уменьшаются еще на один член, что приводит к~сокращению числа комбинаций 


на~8!, так что остается рассмотреть $10! \hm- 8! \hm= 3\,628\,800 \hm- 40\,320 


\hm= 3\,588\,480$ комбинаций. Малоутешительно! 


    


    \smallskip


    


    \textbf{Стратегия 3} (Поиск и~рассуждение). Составим дерево поиска, 


упрощенный вариант которого изображен на рис.~1. 


    


    Ветвь $\mathbf{L} = 1$ $\mathbf{R} = 3$ приводит к~противоречию: 


отрицательному~\textbf{G} в~\mbox{6-м} разряде. По аналогии рассматриваются все остальные 


ветви, причем далее каждая ветвь ветвится еще на один ярус (что не показано 


на рис.~1), поскольку, как указывает Г.~Саймон, <<букве~\textbf{О} не удается 


приписать ка\-кое-ни\-будь единственное значение, удовлетворяющее всем 


полученным условиям. Для этого в~каждом отдельном случае необходимо 


рассмотреть четыре дополнительных варианта. Следовательно, полное дерево 


поиска должно иметь~68~ветвей. Но и~это существенно меньше, чем~10! или 


даже~9!>>~\cite[c.~40]{7-kor}. 


    


 


  











    \smallskip


    


    \textbf{Стратегия 4} (Развитие стратегии поиска и~рассуждения) 


предполагает обнаружение столбцов, относительно которых уже имеется 


достаточная определенность, позволяющая проводить новые подстановки.





\begin{figure} %fig1


\begin{center}


{\small


\begin{tabular}{|ccccccccc|}


\hline


&&&&&&&&\\[-9pt]


$\mathbf{D}=5$ & $\mathbf{T}=0$ & $\mathbf{L}=1$ & $\mathbf{R}=3$& $\mathbf{G}<0$ & &&&\\


&&$\mathbf{L}=2$ & $\mathbf{R}=5$ & $\mathbf{G}=0$ &&&&\\


&&$\mathbf{L}=3$& $\mathbf{R}=7$ & $\mathbf{A}=1$ & $\mathbf{E}=2$ &&&\\


&&&&$\mathbf{A}=2$ & $\mathbf{E}=4$ &&&\\


&&&&$\mathbf{A}=4$ & $\mathbf{E}=8$ &&&\\


&&&&$\mathbf{A}=6$ & $\mathbf{E}=2$ &&&\\


&&&&$\mathbf{A}=8$ & $\mathbf{E}=6$ &&&\\


&&&&$\mathbf{A}=9$ & $\mathbf{E}=8$ &&&\\


&& $\mathbf{L}=4$ & $\mathbf{R}=9$ & $\mathbf{A}=1$& $\mathbf{E}=2$& &&\\


&&&&$\mathbf{A}=2$ & $\mathbf{E}=4$ &&&\\


&&&&$\mathbf{A}=3$ & $\mathbf{E}=6$ &&&\\


&&&&$\mathbf{A}=6$ & $\mathbf{E}=2$ &&&\\


&&&&$\mathbf{A}=7$ & $\mathbf{E}=4$ &&&\\


&&&&$\mathbf{A}=8$ & $\mathbf{E}=6$ &&&\\


&&$\mathbf{L}=6$& $\mathbf{R}=3$ & $\mathbf{G}<0$&&&&\\


&& $\mathbf{L}=7$ & $\mathbf{R}=5$ &&&&&\\


&&$\mathbf{L}=8$& $\mathbf{R}=7$ & $\mathbf{A}=1$ & $\mathbf{E}=3$& &&\\


&&&& $\mathbf{A}=2$ & $\mathbf{E}=5$&&&\\


&&&& $\mathbf{A}=3$ & $\mathbf{E}=7$&&&\\


&&&& $\mathbf{A}=4$ &$\mathbf{E}=9$ & $\mathbf{N}=1$ & $\mathbf{B}=8$ &\\


&&&&  && $\mathbf{N}=2$ & $\mathbf{B}=9$ &\\


&&&&  && $\mathbf{N}=3$ & $\mathbf{G}=0$ &\\


&&&&  && $\mathbf{N}=6$ & $\mathbf{O}=2$ &$\mathbf{G}=1$\\


\hline


\end{tabular}


}


\end{center}


\vspace*{-3pt}





\Caption{Упрощенное дерево поиска для задачи~1~\cite[c.~40]{7-kor}}


\vspace*{-3pt}


\end{figure}


    


    При $\mathbf{D} = 5$ из 1-го разряда следует, что $\mathbf{Т}\hm = 0$. 


Тогда перенос единицы из 1-го во 2-й разряд преобразует сумму 


$\mathbf{L}\hm+ \mathbf{L}\hm = \mathbf{R}$ в~выражение $\mathbf{R}\hm = 


2\mathbf{L}\hm + 1$, где \textbf{R}~--- нечетное число. Кроме того, $\mathbf{R}\hm> 5$, 


поскольку в~6-м разряде имеем соотношение $\mathbf{R}\hm = 5 \hm+ 


\mathbf{G}$. Тогда $\mathbf{R}\hm = 7$ или 9. В~5-м разряде соотношение 


$\mathbf{O}\hm + \mathbf{E}\hm = \mathbf{O}$ справедливо, если 


$\mathbf{E}\hm = 0$ или~9. Но поскольку установлено, что $\mathbf{Т}\hm = 


0$, тогда $\mathbf{Е}\hm = 9$, а $\mathbf{R}\hm = 7$. Таким образом, 


 с~по\-мощью простейших логических рассуждений получаем, что при 


$\mathbf{D}\hm = 5$ и~$\mathbf{T}\hm = 0$, $\mathbf{Е}\hm = 9$ 


и~$\mathbf{R}\hm = 7$.


    


    Поскольку единица перенесена из 2-го разряда в~3-й, то $\mathbf{Е}\hm = 


2\mathbf{А}\hm + 1$, что возможно при $\mathbf{А}\hm = 4$. Но если из 2-го 


разряда в~3-й и~из 1-го во 2-й перенесена единица, то во 2-м разряде имеем 


соотношение не $2\mathbf{L}\hm + 1 \hm= \mathbf{R}$, а~$2\mathbf{L}\hm + 1 


\hm= 10 \hm+ \mathbf{R}$. Тогда $\mathbf{L}\hm = 8$.


    


    Остается распределить цифры~1, 2, 3 и~6 между буквами \textbf{N}, \textbf{B}, \textbf{O} и~\textbf{G}. В 6-м 


разряде имеем $\mathbf{D}\hm + \mathbf{G}\hm + 1\hm = \mathbf{R}$, 


поскольку при любом значении~O в~6-й разряд переносится единица из 5-го 


разряда. Тогда $\mathbf{G}\hm = 1$, а~для N, B и~O имеем ряд цифр~2, 3 и~6, 


т.\,е.\ осталось перебрать $3! \hm= 6$ различных возможностей. Методом 


<<проб и~ошибок>> получаем $\mathbf{N}\hm = 6$, $\mathbf{B}\hm = 3$ 


и~$\mathbf{O}\hm = 2$. 





\section{Альтернативные способы оптимизации поиска при~решении~криптоарифметических задач}





%\vspace*{-3pt}





    Приведем подходы к~решению задач Саймона, отличные от способов 


самого автора. При этом используем логические рассуждения, основанные на 


блок-схе\-мах алгоритма, причем оптимизацию блок-схем за счет сужения 


области поиска будем проводить с~использованием следующих опор: 


    \begin{enumerate}[(1)]


    \item проверки четности-нечетности замещающих чисел (цифр);


    \item наличия свободного ряда цифр;


    \item наличия или отсутствия переноса единицы из предыдущего разряда 


в~последующий, что непосредственно влияет на соотношение между искомыми 


числовыми значениями. 


    \end{enumerate}


    


    Если перенос единицы осуществляется {\bfseries\textit{в}}~разряд, то в~формуле $X \hm+ Y 


\hm= Z$ к~левой части добавляется единица: $X \hm+ 1 \hm+ Y \hm= Z$. Если 


перенос единицы осуществляется {\bfseries\textit{из}} разряда, то в~формуле $X \hm+ Y \hm= Z$ 


к~правой части добавляется число десять: $X \hm+ Y \hm= 10 


\hm+ Z$. Здесь $X$, $Y$ и~$Z$~--- цифры одного разряда при записи в~столбик.


    


   \smallskip


    


    \textbf{Решение задачи~1.}\ Имеем 10~букв: \textbf{D}, 


    \textbf{O}, \textbf{N}, \textbf{A}, \textbf{L}, \textbf{G}, 


    \textbf{E}, \textbf{R}, \textbf{B} и~\textbf{T} и~10~цифр: 0, 1, 2, 3, 4, 5, 6, 7, 8 и~9. При этом формальный перечень возможных 


комбинаций имеет~16~ветвей (цепочек), на основании анализа которых можно 


сократить число переборов, оптимизировав таким образом решение задачи.


    


    Из 2-го разряда получаем, что $\mathbf{L}\hm= 9\hm+\mathbf{R}$ или 


$\mathbf{L}\hm= (\mathbf{R}\hm-1)/2$; следовательно, $\mathbf{R}$~--- число нечетное. В 5-м 


разряде имеет смысл единственное соотношение: $\mathbf{О}\hm + 1 \hm+ 


\mathbf{Е}\hm = 10 \hm+ \mathbf{О}$, откуда следует, что $\mathbf{Е}\hm = 9$. 


Тогда из общей схемы комбинаций можно удалить \textbf{12~цепочек} (по три 


для каждой из четырех ветвей в~3-м разряде), сократив число комбинаций до 


\textbf{четырех} и~пронумеровав их римскими цифрами~I--IV (рис.~2). На 


рис.~2 в~скобках указан номер разряда от~(1) до~(6); кружок со стрелкой~--- 


это перенос единицы из предыдущего разряда в~последующий, 


а~перечеркнутый кружок символизирует отсутствие переноса. 


    


    


 


 


    Дальнейшая оптимизация поиска осуществляется в~3-м разряде для 


комбинаций (см.\ рис.~2): 


    \begin{enumerate}[I]


    \item  $\mathbf{A}=(9+\mathbf{E})/2=9$;


     \item $\mathbf{A}=(\mathbf{E}-1)/2=4$;


     \item $\mathbf{A}=(10+\mathbf{E})/2=9{,}5$;


     \item $\mathbf{A}=\mathbf{E}/2=4{,}5$.


     \end{enumerate}


    


    Поскольку~\textbf{A} должно быть целым числом и~к~тому же отличным от~9, то 


комбинации~I, III и~IV являются тупиковыми, т.\,е.\ остается рассмотреть всего 


\textbf{одну комбинацию}~--- под номером~II, для которой имеем: 


$\mathbf{D}\hm = 5$, $\mathbf{T}\hm = 0$ и~$\mathbf{E}\hm = 9$. Дальнейшие 


рассуждения сведены в~табл.~1, демонстрирующую поисковые действия при 


движении слева направо (в~верхней строке в~скобках указан номер разряда из 


рис.~2). 


 


 


 


 \noindent


 \textbf{Ответ}: $\mathbf{D} = 5$; $\mathbf{T}\hm = 0$; $\mathbf{E}\hm = 9$; 


$\mathbf{A}\hm = 4$; $\mathbf{G}\hm = 1$; $\mathbf{R}\hm = 7$; $\mathbf{L} \hm= 8$; 


$\mathbf{N}\hm = 6$; $\mathbf{В}\hm = 3$; $\mathbf{O}\hm = 2$.





 


 


 \noindent


 \textbf{Проверка}:\ \ \ \ \ \ 526\,485\\[-19pt]


 


 \noindent


 \hphantom{ \textbf{Проверк}\ }\ \ \ $+$\\[-19pt]


 


 \noindent 


\hphantom{ \textbf{Проверк\,}:\ }\ \ \ \ \ 197\,485


 


 \noindent


\hphantom{ \textbf{Проверк\,}:\ }\ \ \ \ \ $\overline{723\,970}$


    


    Таким образом, верный ответ получен за четыре итерации при 


рассмотрении одной комбинации~--- под номером~II. Оставшейся цифрой~2 


можно заместить букву \textbf{О}, что никак не повлияет на конечный результат. 





\begin{figure} %fig2


 \vspace*{1pt}


  \begin{center}


 \mbox{%


 \epsfxsize=125.589mm 


\epsfbox{kor-2.eps}


 }


\end{center}


\vspace*{-9pt}


 \Caption{Промежуточный этап оптимизации поиска при решении задачи~1 (в верхней 


строке в~скобках указан номер разряда, а в~столбике слева~--- номер комбинации)}


% \end{figure}


%\vspace*{-6pt}


 \renewcommand{\figurename}{\protect\bf Таблица}


\setcounter{figure}{0}


% \begin{table}


{\footnotesize %tabl1


 \begin{center}


 \Caption{Заключительные итерации в~решении задачи~1}


 \vspace*{2ex}


 


 \tabcolsep=2pt


 \begin{tabular}{|c|c|c|c|c|c|c|c|c|}


 \hline


 № & (3) & & (6) & & (2)& & (4)& (5)\\


\tabcolsep=0pt\begin{tabular}{c}ком-\\ бина-\\ ции\end{tabular}&


\tabcolsep=0pt\begin{tabular}{c}\textbf{А}\end{tabular}&


\tabcolsep=0pt\begin{tabular}{c}Свободный\\ ряд\end{tabular}&\tabcolsep=0pt\begin{tabular}{c} $\mathbf{G} + 6 = \mathbf{R}$,\\ 


\textbf{R} и~\textbf{G}~---\\ нечетные числа\end{tabular}&


\tabcolsep=0pt\begin{tabular}{c}Свободный\\ ряд\end{tabular}&\tabcolsep=0pt\begin{tabular}{c} $\mathbf{L}=9+\mathbf{R}$\end{tabular} &


\tabcolsep=0pt\begin{tabular}{c}Свободный\\ ряд\end{tabular}&\tabcolsep=0pt\begin{tabular}{c} $\mathbf{N} = 3 + \mathbf{B}$\end{tabular} 


&\tabcolsep=0pt\begin{tabular}{c} \textbf{O}  \end{tabular}\\


\hline


 II&4&1, 2, 3, 6, 7, 8&$\mathbf{G} = 1$; $\mathbf{R = 7}$&2, 3, 6, 8&$\mathbf{L} = 8$


 &2, 3, 6&\tabcolsep=0pt\begin{tabular}{c}$\mathbf{N} = 6$\\ $\mathbf{B} = 3$\end{tabular}&2\\


 \hline


 \end{tabular}


 \end{center}


 }


 \vspace*{-6pt}


% \end{table}


\end{figure}


 


 


 \renewcommand{\figurename}{\protect\bf Рис.}


\renewcommand{\tablename}{\protect\bf Таблица}


\setcounter{table}{1}





    


    \smallskip


    


    \textbf{Решение задачи~2.}\ Имеем 9 букв: \textbf{C}, 


    \textbf{R}, \textbf{O}, \textbf{S}, \textbf{A}, \textbf{D}, \textbf{N}, \textbf{G} и~\textbf{E} и~10 


цифр: 0, 1, 2, 3, 4, 5, 6, 7, 8 и~9, причем очевидно, что $\mathbf{D} \hm= 1$, поскольку 


происходит перенос разряда из 5-го в~6-й. При этом формальный перечень 


возможных комбинаций включает~16~ветвей. Показательно, что сам Герберт 


Саймон, рассматривая оптимизацию поиска в~задаче~1, отмечал, что в~задаче~2, решение которой сам он не приводит, не удается <<столь же 


радикально сократить поиск>>~\cite[с.~41]{7-kor}. Правота его слов будет 


подтверждена в~ходе дальнейших рас\-суж\-де\-ний. 


{\looseness=1





}


    


    {В 1-м разряде} имеем два начальных соотношения: (1)~$ 2\mathbf{S}\hm = 10 \hm+ \mathbf{R}$ и~(2)~$ 2\mathbf{S}\hm = \mathbf{R}$, 


где число $\mathbf{R}$~--- четное для обоих случаев. Если построить 


развернутую блок-схему полного набора комбинаций, исходящих из двух 


начальных соотношений, то она покажет, что каждое начальное соотношение 


разворачивается в~восемь цепочек. В~целом это соответствует~16~поисковым 


комбинациям (пока без учета возможных сочетаний цифр, которых в~пределах 


одной цепочки может быть несколько). 


    


    {Для начального соотношения} $2\mathbf{S}\hm = 


\mathbf{10}\hm + \mathbf{R}$ число~$\mathbf{R}$ может принимать значения 


0, 2, 4, 6 и~8. Тогда $\mathbf{S}\hm = 5$, 6, 7, 8 и~9. Для 2-го разряда этого 


соотношения существуют два варианта: $\mathbf{S}\hm + \mathbf{D}\hm + 1\hm 


= 10 \hm+ \mathbf{Е}$ и~$\mathbf{S}\hm + \mathbf{D}\hm + 1 \hm= \mathbf{E}$, 


что при $\mathbf{D}\hm = 1$ приводит к~простым формулам: $\mathbf{S}\hm = 8 \hm+ 


\mathbf{E}$ и~$\mathbf{S}\hm = \mathbf{E}\hm- 2$ соответственно. 


    


    \textit{В первом варианте} 2-го разряда соотношение $\mathbf{S}\hm = 8 


\hm+ \mathbf{E}$ может быть выполнено только при $\mathbf{Е}\hm = 0$. 


Тогда $\mathbf{S}\hm = 8$, а~$\mathbf{R}\hm = 6$ (при этом образуется 


свободный ряд цифр~2, 3, 4, 5, 7 и~9). 


    


    \textit{Во втором варианте} 2-го разряда для $\mathbf{S}\hm = 


\mathbf{E}\hm- 2$ возможны три сочетания цифр для \textbf{R}, \textbf{S} и~\textbf{E}: $\mathbf{R}\hm 


= 0$, $\mathbf{S} \hm= 5$, $\mathbf{E}\hm = 7$ (свободный ряд цифр 2, 3, 4, 6, 8 и~9); 


$\mathbf{R}\hm = 2$, $\mathbf{S}\hm = 6$ и~$\mathbf{E}\hm = 8$ (свободный ряд 


цифр~0, 3, 4, 5, 7 и~9); $\mathbf{R}\hm = 4$, $\mathbf{S}\hm = 7$ 


и~$\mathbf{E}\hm = 9$ (свободный ряд цифр~0, 2, 3, 5, 6 и~8). 


    


    Следовательно, {для первого начального соотношения} 


$2\mathbf{S}\hm = {10}\hm + \mathbf{R}$ на этапе \mbox{2-го} разряда 


открывается возможность продолжить поиск по всем восьми цепочкам: 


четырем для соотношения $\mathbf{S}\hm = 8 \hm+ \mathbf{E}$ при 


единственном сочетании цифр: $\mathbf{D}\hm = 1$, $\mathbf{E}\hm = 0$, 


$\mathbf{S}\hm = 8$ и~$\mathbf{R}\hm = 6$ и~свободном ряде цифр (2, 3, 4, 5, 7 и~9) 


и~четырем для соотношения $\mathbf{S}\hm = \mathbf{E}\hm- 2$ при трех 


вариантах сочетаний цифр:


    \begin{enumerate}[(1)]


    \item $\mathbf{D} = 1$, $\mathbf{E} = 7$, $\mathbf{S} = 5$ и~$\mathbf{R}\hm 


= 0$ и~свободном ряде цифр (2, 3, 4, 6, 8 и~9);


\item $\mathbf{D} = 1$, $\mathbf{E} = 8$, $\mathbf{S} = 6$ 


и~$\mathbf{R}\hm = 2$ и~свободном ряде цифр (0, 3, 4, 5, 7 и~9);


\item $\mathbf{D} = 1$, $\mathbf{E} = 9$, $\mathbf{S}\hm = 7$ 


и~$\mathbf{R}\hm = 4$ и~свободном ряде цифр (0, 2, 3, 5, 6 и~8).


\end{enumerate}


    


    Таким образом, {для первого начального соотношения} 


$2\mathbf{S}\hm = \mathbf{10}\hm + \mathbf{R}$ в~пределах восьми цепочек 


имеем $4\times 1 \hm+ 4\times3 \hm= {16}$ {разных вариантов}, 


среди которых будем искать решение. 


    


    {Для второго начального соотношения} $2\mathbf{S}\hm = 


\mathbf{R}$ число~\textbf{R} может принимать значения~4, 6 или~8. Тогда $\mathbf{S}\hm 


= 2$, 3 или~4. Для 2-го разряда этого соотношения могут рассматриваться два 


варианта: $\mathbf{S}\hm + \mathbf{D}\hm = 10 \hm+ \mathbf{Е}$ 


и~$\mathbf{S}\hm + \mathbf{D}\hm = \mathbf{E}$, что при $\mathbf{D}\hm = 1$ 


приводит к~простым формулам: $\mathbf{S} \hm= 9 \hm+ \mathbf{E}$ и~$\mathbf{S}\hm 


= \mathbf{E}\hm- 1$ соответственно. 


    


    \textit{В первом варианте} 2-го разряда соотношение $\mathbf{S}\hm = 9 


\hm+ \mathbf{E}$ не имеет смысла при всех сочетаниях значений~\textbf{R} и~\textbf{S}, 


поэтому все четыре цепочки комбинации $(2\mathbf{S}\hm = \mathbf{R})\hm\to 


(\mathbf{S}\hm = 9 \hm+ \mathbf{E})$ можно не рассматривать, поскольку они 


тупиковые. Так исключаем из рассмотрения сразу четыре цепочки. 


    


    \textit{Во втором варианте} 2-го разряда для $\mathbf{S}\hm = 


\mathbf{E}\hm- 1$ возможны три сочетания цифр для \textbf{R}, \textbf{S} и~\textbf{E}: $\mathbf{R}\hm 


= 4$, $\mathbf{S}\hm = 2$ и~$\mathbf{E}\hm = 3$ (свободный ряд~0, 5, 6, 7, 8 и~9); 


$\mathbf{R}\hm = 6$, $\mathbf{S}\hm = 3$ и~$\mathbf{E}\hm = 4$ (свободный ряд 


0, 2, 5, 7, 8 и~9); $\mathbf{R}\hm = 8$, $\mathbf{S}\hm = 4$ и~$\mathbf{E}\hm = 5$ 


(свободный ряд 0, 2, 3, 6, 7 и~9). 


    


    Следовательно, {для второго начального соотношения} 


$2\mathbf{S}\hm= \mathbf{R}$ можно продолжать поиск по четырем 


комбинациям из восьми, соответствующим соотношению $\mathbf{S}\hm = 


\mathbf{E}\hm- 1$ на этапе 2-го разряда, в~каждой цепочке которого имеем по 


три сочетания цифр~--- итого $4\times 3 \hm= {12}$~{разных 


вариантов}, среди которых, возможно, скрывается решение задачи. 


    


    Итак, после предварительных рассуждений вся задача разбивается на три 


независимые подзадачи по четыре комбинации в~каждой, поскольку общее 


число комбинаций, изначально равное~16, сократилось до~12 за счет 


исключения четырех комбинаций с~начальным соотношением $(2\mathbf{S}\hm 


= \mathbf{R})\hm\to  (\mathbf{S}\hm = 9 \hm+ \mathbf{E})$. Остаются 


4~комбинации для $(2\mathbf{S}\hm = 10 \hm+ \mathbf{R}) \hm\to 


(\mathbf{S}\hm = 8 \hm+ \mathbf{E})$, а еще~8~комбинаций~--- четыре для 


$(2\mathbf{S}\hm = 10 \hm+ \mathbf{R})\hm\to (\mathbf{S}\hm = \mathbf{E}\hm- 


2)$ и~четыре для $(2\mathbf{S}\hm = \mathbf{R}) \hm\to (\mathbf{S}\hm = 


\mathbf{E}\hm- 1)$~--- <<обогащаются>> тремя сочетаниями цифр каждая. 


Следовательно, всего необходимо проверить $4\times1 \hm+ 4\times3 \hm+ 


4\times3 \hm= 28$ вариантов возможных соотношений между буквами 


и~цифрами. Это доказывает правоту Герберта Саймона о том, что решение 


задачи~2 требует гораздо б\'{о}льших усилий за счет невозможности 


радикального сокращения области поиска (напомним, что в~задаче~1 


из~16~комбинаций были исключены~12). 


    


    {Рассмотрим кратко ход решения задачи~2.} Как было показано 


выше, для начального соотношения $(2\mathbf{S}\hm = 10 \hm+ \mathbf{R}) 


\hm\to  (\mathbf{S}\hm = 8 \hm+ \mathbf{E})$ имеется единственный возможный 


вариант: $\mathbf{R} \hm= 6$, $\mathbf{S} \hm= 8$ и~$\mathbf{E}\hm= 0$ и~свободный ряд цифр: 2, 


3, 4, 5, 7 и~9. Далее воспользуемся следующими формулами для вычисления 


остальных бук\-вен\-но-чис\-ло\-вых соответствий в~четырех соответствующих 


цепочках: 


    \begin{enumerate}[(1)]


    \item  $\mathbf{O}\hm = 9 \hm+ \mathbf{N}\hm- \mathbf{R}\hm = 


\mathbf{N}\hm + 3$; $\mathbf{C}-\mathbf{A}\hm = 9 \hm- \mathbf{R}\hm = 3$; 


$\mathbf{G}\hm = \mathbf{O} \hm+ \mathbf{A}\hm- 9$; 


    \item  $\mathbf{O} = \mathbf{N}- \mathbf{R}- 1 \hm= \mathbf{N}- 7$; 


$\mathbf{C} -\mathbf{A}\hm= 10 \hm- \mathbf{R}\hm = 4$; $\mathbf{G} \hm= \mathbf{O}\hm 


+ \mathbf{A} \hm- 9$; 


    \item  $\mathbf{O} = 10 + \mathbf{N}- \mathbf{R}\hm = \mathbf{N}\hm + 4$; 


$\mathbf{C}\hm- \mathbf{A}\hm = 9 \hm- \mathbf{R} \hm= 3$; $\mathbf{G}\hm = 


\mathbf{O}\hm + \mathbf{A}\hm + 1$; 


    \item $\mathbf{O} = \mathbf{N}- \mathbf{R}\hm = \mathbf{N}\hm- 6$; 


$\mathbf{C}\hm- \mathbf{A}\hm= 10 \hm- \mathbf{R} \hm= 4$. 


    \end{enumerate}


    


    Расчеты показывают, что ни одна из рассмотренных цепочек не приводит к~приемлемому результату. 


    


    Напомним, что во всех четырех цепочках для комбинации 


$(2\mathbf{S}\hm = 10 \hm+ \mathbf{R}) \hm\to (\mathbf{S}\hm = 


\mathbf{E}\hm- 2)$ имеются по три набора значений: $\mathbf{R}\hm = 0$, 2, 4; 


$\mathbf{S}\hm = 5$, 6, 7; $\mathbf{E}\hm = 7$, 8, 9. Во всех четырех цепочках 


для комбинации $(2\mathbf{S}\hm = \mathbf{R}) \hm\to (\mathbf{S}\hm = 


\mathbf{E}\hm- 1)$ также имеется по три набора значений: $\mathbf{R}\hm = 4$, 


6, 8; $\mathbf{S}\hm = 2$, 3, 4; $\mathbf{E}\hm = 3$, 4, 5. Проведя соответствующие 


вычисления, получаем, что к~приемлемому результату приводит только одна из 


цепочек для $(2\mathbf{S}\hm = \mathbf{R}) \hm\to (\mathbf{S}\hm = 


\mathbf{E}\hm- 1)$, соответствующая значениям: $\mathbf{R}\hm = 6$, 


$\mathbf{S}\hm = 3$ и~$\mathbf{E}\hm = 4$ (табл.~2).


 


 \begin{table}\small %tabl2


 \begin{center}


 \Caption{Фрагмент решения задачи~2 для комбинации $(2\mathbf{S}\hm = 


\mathbf{R}) \hm\to  (\mathbf{S}\hm = \mathbf{E}\hm- 1)$}


 \vspace*{2ex}


 


 \tabcolsep=7pt


 \begin{tabular}{|c|c|c|c|c|}


 \hline


 \multicolumn{1}{|c|}{\raisebox{-18pt}[0pt][0pt]{Шаги 1--2}}&\textbf{R}&4&6&8\\


% \cline{2-5}


 &\textbf{S}&2&3&4\\


 % \cline{2-5}


 &\textbf{E}&3&4&5\\


%  \cline{2-5}


 &{Свободный ряд}&0~5~6~7~8~9&0~2~5~7~8~9&0\ 2\ 3\ 6\ 7\ 9\\


 \hline


 &&\multicolumn{3}{c|}{\ }\\[-9pt]


&\textbf{О}&\multicolumn{3}{c|}{$\mathbf{N} - \mathbf{R}$}\\


  \cline{3-5}


 &&\multicolumn{3}{c|}{\ }\\[-9pt]


  &\textbf{O}&$\mathbf{N} - 4$ &$\mathbf{N} - 6$ &$\mathbf{N}- 8$ \\


  % \cline{2-5}


  Шаг 3&\textbf{N}&9&8&Нет решения\\


 % \cline{2-5}


 &\textbf{О}&5&2&Нет решения\\


 % \cline{2-5}


 &{Свободный ряд}&0 6 7 8&0 5 7 9 &---\\


 \hline


 &&\multicolumn{3}{c|}{\ }\\[-10pt]


&$\mathbf{С} - \mathbf{A}$&\multicolumn{3}{c|}{$10 - \mathbf{R}$}\\


  \cline{3-5}


 &$\mathbf{C}-\mathbf{A}$&6&4&---\\


 % \cline{2-5}


  Шаг 4&\textbf{А}&0&5&---\\


 % \cline{2-5}


  &\textbf{С}&6&9&---\\


 %  \cline{2-5}


 &{Свободный ряд}&7~8&0~7&---\\


 \hline


 &&\multicolumn{3}{c|}{\ }\\[-10pt]


\multicolumn{1}{|c|}{\raisebox{-6pt}[0pt][0pt]{Шаг 5}}&\textbf{G}&\multicolumn{3}{c|}{$\mathbf{O} + \mathbf{A}$ }\\


 \cline{3-5}


  &&&&\\[-9pt]


 &\textbf{G}&Нет решения&\textbf{РЕШЕНИЕ: 7}&---\\


 \hline


 \end{tabular}


 \end{center}


 \vspace*{-6pt}


 \end{table}


 


 \noindent


 \textbf{Ответ}: $\mathbf{C} = 9$, $\mathbf{R} = 6$, $\mathbf{O}\hm = 2$, 


$\mathbf{S}\hm = 3$, $\mathbf{A}\hm = 5$, $\mathbf{D}\hm = 1$, 


$\mathbf{N}\hm = 8$, $\mathbf{G}\hm = 7$, $\mathbf{E}\hm = 4$. Цифра~0 


осталась невостребованной. 





\pagebreak


 


 \noindent


 \textbf{Проверка}:\ \ \ \ \ \ 96\,233\\[-19pt]


 


\noindent 


 \hphantom{ \textbf{Проверк}\ }\ \ \ $+$\\[-19pt]


               


               \noindent


             \hphantom{\textbf{Проверк\,}:\ } \ \ \ \ \ 62\,513


               


               \noindent


             \hphantom{\textbf{Провер\,}:\ }\ \ \ \ \ \ $\overline{158\,746}$


    


    Итак, для получения ответа пришлось перебрать~12~комбинаций 


возможных решений, рассмотрев в~целом 28 переборов (итераций): 4~перебора 


для комбинации $(2\mathbf{S}\hm = 10 \hm+ \mathbf{R}) \hm\to 


(\mathbf{S}\hm = 8 \hm+ \mathbf{E})$ и~по~12~переборов для комбинаций 


$(2\mathbf{S}\hm = 10 \hm+ \mathbf{R}) \hm\to  (\mathbf{S}\hm = 


\mathbf{E}\hm- 2)$ и~$(2\mathbf{S}\hm = \mathbf{R}) \hm\to (\mathbf{S}\hm = 


\mathbf{E}\hm- 1)$. Если сравнивать число переборов, необходимых для 


решения задач~1 и~2, то получаем, что в~задаче~1 проделаны~4~итерации для 


одной комбинации, а~в~задаче~2~--- 28~итераций для~12~комбинаций. 


    


    Отметим, что Герберт Саймон даже в~наиболее продвинутой стратегии 


оптимизации поиска~--- под номером~4~--- говорит о~методе <<проб 


и~ошибок>>, к~которому он прибегает на заключительном этапе оптимизации,


тогда как в~предложенной методике оптимизации поиск на любом шаге 


итерации не ориентируется на метод <<проб и~ошибок>>, а~носит 


обоснованный и~целенаправленный характер, опирающийся на четкие 


арифметические правила и~закономерности. 





\vspace*{-6pt}





\section{Заключение}








\vspace*{-4pt}





    Оптимизация поиска при решении криптоарифметических задач служит 


эффективным инструментом развития логического мышления школьников, 


способствующего формированию алгоритмических навыков как необходимых 


компонентов инженерной культуры. При соответствующем методическом 


\mbox{сопровождении} сложные переборные задачи, которые в~настоящее время 


встречаются преимущественно на предметных олимпиадах, могут быть 


включены в~программу углубленного курса информатики на уровне основного 


общего образования. Впоследствии это может послужить дополнительным 


стимулом для освоения учащимися методов динамического 


программирования~--- эффективной технологии решения многих классов 


переборных задач. 


    


 {\small\frenchspacing


{%\baselineskip=10.8pt


\addcontentsline{toc}{section}{Литература}


\begin{thebibliography}{9}


\bibitem{1-kor}


\Au{Корчажкина О.\,М.} Криптоарифметические и~другие переборные задачи на уроках 


информатики~/\!\!/ Актуальные проблемы методики обучения информатике и~математике 


в~современной школе: мат-лы Междунар. научн.-практич. ин\-тер\-нет-кон\-фе\-рен\-ции.~--- 


М.: МПГУ, 2022. {\sf http://news.scienceland.ru/2022/04/17/криптоарифметические-и-другие-переб/}. 


\bibitem{2-kor}


\Au{Булычёв В.\,А.} Методы программирования: переборные алгоритмы. {\sf 


https://algolist. manual.ru/maths/combinat}.





\bibitem{5-kor} %3


\Au{Окулов С.\,М., Пестов~О.\,А.} Динамическое программирование.~--- М.: БИНОМ, 


 2015. 299~с.


\bibitem{3-kor} %4


\Au{Довгалюк П.\,М.} Динамическое программирование и~все-все-все: Как решать 


олимпиадные и~<<жизненные>> программистские задачи.~--- М.: ЛЕНАНД, 2021. 200~с. 


\bibitem{4-kor} %5


\Au{Корчажкина О.\,М.} Технология динамического программирования как инструмент 


развития инженерного мышления старшеклассников~/\!\!/ Преподавание информационных 


технологий в~Российской Федерации: Сб. научных тр. 19-й Открытой Всеросс. 


конф.~--- М.: 1С-Паб\-ли\-шинг, 2021. С.~86--88.





\bibitem{6-kor}


\Au{Прокопцев А.\,А.} Методические подходы к~изучению основ динамического 


программирования на уровне среднего общего образования~/\!\!/ Актуальные проб\-ле\-мы 


методики обучения информатике и~математике в~современной школе: мат-лы Междунар. 


научн.-практич. ин\-тер\-нет-кон\-фе\-рен\-ции.~--- М.: МПГУ, 2022. {\sf 


http://news. scienceland.ru/2022/04/17/методические-подходы-к-изучению-осно/}.


\bibitem{7-kor}


\Au{Саймон Г.} Науки об искусственном~/ Пер. с~англ.~--- 2-е изд.~--- М.: Едиториал УРСС, 


2004. 144~с. 


(\Au{Simon~H.\,A.} The sciences of the artificial.~--- 3rd ed.~--- Cambridge, MA, USA: 


MIT Press, 1996. 248~p.)





\bibitem{8-kor}


\Au{Пенроуз Р.} Новый ум короля: О компьютерах, мышлении и~законах физики~/ Пер. 


с~англ.~--- 3-е изд.~--- М.: ЛКИ, 2008. 400~с.


(\Au{Penrose~R.} The Emperor's new mind: Concerning computers, minds and the laws of 


physics.~--- Oxford, U.K.: Oxford University Press, 2002. 640~p.)





 \end{thebibliography}


} }











\vspace*{-6pt}





\hfill{\small\textit{Поступила в~редакцию 15.09.22}}





%\vspace*{8pt}





%\hrule





%\vspace*{2pt}





\newpage





\vspace*{-28pt}





%\hrule





%\vspace*{6pt}





\def\tit{SEARCH OPTIMIZATION WHILE SOLVING ENUMERATION PROBLEMS 


IN~AN~ADVANCED COMPUTER SCIENCE COURSE AT~THE~LEVEL 


OF~BASIC GENERAL EDUCATION}











\def\titkol{Search optimization while solving enumeration problems 


in~an~advanced computer %science 


course} % at~the~level of~basic general education}





\def\aut{O.\,M.~Korchazhkina}





\def\autkol{O.\,M.~Korchazhkina}





\titel{\tit}{\aut}{\autkol}{\titkol}





\vspace*{-8pt}





\noindent


Federal Research Center ``Computer Science and Control'' of the Russian Academy of Sciences, 


44-2~Vavilov Str., Moscow 119333, Russian Federation 





\def\leftfootline{\small{\textbf{\thepage}


\hfill Sistemy i Sredstva Informatiki~---


Systems and Means of Informatics 2022 vol~32 no\,4}


}%


 \def\rightfootline{\small{Sistemy i Sredstva Informatiki~---


 Systems and Means of Informatics 2022 vol~32 no\,4


\hfill \textbf{\thepage}}}





\vspace*{3pt}











\Abste{The article discusses a few ways to activate high school students' logical 


thinking to form their algorithmic skills in computer studies lessons. The 


research 


analyses the methods for optimizing in traditional cryptoarithmetic problems that 


make up a subclass of enumeration tasks which were proposed by Herbert Simon to 


illustrate the ways of mindset when it is necessary to carry out numerous iterations of 


possible combinations. A~search optimization visual


method is proposed for solving 


two Simon problems based on the patterns of


digit transfer during summation and 


presented in the form of simple algebraic


relations. The method can be easily 


algorithmized and implemented with dynamic programming technology to optimize 


the solution of iterative tasks.}





\KWE{optimization of a search; enumeration\!/\!iterative problem\!/\!task;  


crypto-arithmetic problem; dynamic programming}








\DOI{10.14357\!/\!08696527220414}  





%\vspace*{-10pt}








%\Ack


%\noindent











\vspace*{-9pt}





\renewcommand{\bibname}{\protect\rmfamily References}


%\renewcommand{\bibname}{\large\protect\rm References}





{\small\frenchspacing


{ %\baselineskip=12pt


\addcontentsline{toc}{section}{References}


\begin{thebibliography}{9}





\vspace*{-2pt}








\bibitem{1-kor-1}


\Aue{Korchazhkina, O.\,M.} 2022. Kriptoarifmeticheskie i~drugie perebornye 


zadachi na urokakh informatiki [Cryptoarithmetic and other iterative problems in 


computer science lessons]. \textit{Scientific and Practical Conference 


(International) ``Actual Problems of Teaching Methods of Computer Science and 


Mathematics in a Modern School'' Proceedings}. Moscow: MPGU. Available at: {\sf 


http://news.scienceland.ru/2022/04/17/ криптоарифметические-и-другие-переб/} 


(accessed October~10, 2022).


\bibitem{2-kor-1}


\Aue{Bulychev, V.\,A.} Metody programmirovaniya: perebornye algoritmy 


[Programming methods: Enumeration algorithms]. Available at: {\sf 


https://algolist.manual.ru/maths/\linebreak combinat/} (accessed  October~10, 2022).





\bibitem{5-kor-1} %3


\Aue{Okulov, S.\,M., and O.\,A.~Pestov}. 2015. \textit{Dinamicheskoe 


programmirovanie} [Dynamic programming]. Moscow: BINOM. 299~p.


\bibitem{3-kor-1} %4


\Aue{Dovgalyuk, P.\,M.} 2021. \textit{Dinamicheskoe programmirovanie  


i~vse-vse-vse: Kak reshat' olim\-pi\-ad\-nye i~``zhiznennye'' programmistskie zadachi} 


[Dynamic programming and all-all-all: How to solve Olympiad and ``life'' 


programming tasks]. Moscow: LENAND. 200~p.


\bibitem{4-kor-1} %5


\Aue{Korchazhkina, O.\,M.} 2021. Tekhnologiya dinamicheskogo 


programmirovaniya kak instrument razvitiya inzhenernogo myshleniya 


starsheklassnikov [Dynamic programming technology as a~tool to develop high 


school students' engineering mindsets]. \textit{Prepodavanie informatsionnykh tekhnologiy v~Rossiyskoy 


Federatsii: Sb. nauchnykh tr. 19-y  Otkrytoy Vseross. konf.} [19th All-Russian Open Conference ``Teaching Information Technology 


in the Russian Federation'' Proceedings]. Moscow: 1C-Publishing. 86--88.





\bibitem{6-kor-1}


\Aue{Prokoptsev, A.\,A.} 2022. Metodicheskie podkhody k~izucheniyu osnov 


di\-na\-mi\-che\-sko\-go programmirovaniya na urovne srednego obshchego obrazovaniya 


[Methodological approaches to the study of the basics of dynamic programming at 


the level of secondary general education]. \textit{Scientific and Practical Conference 


(International)``Actual Problems of Teaching Methods of 


Computer Science and Mathematics in a~Modern School'' Proceedings}. Moscow: 


MPGU. Available at: {\sf  


http://news.scienceland.ru/\linebreak  2022/04/17/методические-подходы-к-изучению-осно/} 


(accessed October~10, 2022).


\bibitem{7-kor-1}


\Aue{Simon, H.\,A.} 2019. \textit{The sciences of the artificial}. 3rd ed. Cambridge, 


MA: MIT Press. 231~p.


\bibitem{8-kor-1}


\Aue{Penrose, R.} 2016. \textit{The Emperor's new mind: Concerning computers, 


minds and the laws of physics}. Oxford, U.K.: Oxford University Press. 640~p.


\end{thebibliography}


} }





\vspace*{-3pt}





\hfill{\small\textit{Received September 15, 2022}} 





\vspace*{-12pt} 





\Contrl





\vspace*{-4pt}





\noindent


\textbf{Korchazhkina Olga M.} (b.\ 1953)~--- Candidate of Science (PhD) in 


technology, senior scientist, Institute of Cybernetics and Educational 


Computing of Federal Research Center ``Computer Science and Control'' of the 


Russian Academy of Sciences, 44-2~Vavilov Str., Moscow 119333, Russian 


Federation; \mbox{olgakomax@gmail.com}





\label{end\stat}





\renewcommand{\bibname}{\protect\rm Литература} 


